\documentclass[10pt,a4paper]{amsart}
\usepackage{amsmath,amssymb,amsthm,graphicx}
\usepackage{verbatim}
\newenvironment{reference}{\comment}{\endcomment}
%\newenvironment{reference}{}{}
\newenvironment{slogan}{\comment}{\endcomment}
\newenvironment{history}{\comment}{\endcomment}
% For commutative diagrams we use Xy-pic
\usepackage[all]{xy}

% We use 2cell for 2-commutative diagrams.
\xyoption{2cell}
% Theorem environments.
%
\theoremstyle{plain}
\newtheorem{theorem}[subsection]{Theorem}
\newtheorem{proposition}[subsection]{Proposition}
\newtheorem{lemma}[subsection]{Lemma}

\theoremstyle{definition}
\newtheorem{definition}[subsection]{Definition}
\newtheorem{example}[subsection]{Example}
\newtheorem{exercise}[subsection]{Exercise}
\newtheorem{situation}[subsection]{Situation}

\theoremstyle{remark}
\newtheorem{remark}[subsection]{Remark}
\newtheorem{remarks}[subsection]{Remarks}

\numberwithin{equation}{subsection}

% Macros
%
\def\lim{\mathop{\mathrm{lim}}\nolimits}
\def\colim{\mathop{\mathrm{colim}}\nolimits}
\def\Spec{\mathop{\mathrm{Spec}}}
\def\Hom{\mathop{\mathrm{Hom}}\nolimits}
\def\Ext{\mathop{\mathrm{Ext}}\nolimits}
\def\SheafHom{\mathop{\mathcal{H}\!\mathit{om}}\nolimits}
\def\SheafExt{\mathop{\mathcal{E}\!\mathit{xt}}\nolimits}
\def\Sch{\mathit{Sch}}
\def\Mor{\mathop{\mathrm{Mor}}\nolimits}
\def\Ob{\mathop{\mathrm{Ob}}\nolimits}
\def\Sh{\mathop{\mathit{Sh}}\nolimits}
\def\NL{\mathop{N\!L}\nolimits}
\def\CH{\mathop{\mathrm{CH}}\nolimits}
\def\proetale{{pro\text{-}\acute{e}tale}}
\def\etale{{\acute{e}tale}}
\def\QCoh{\mathit{QCoh}}
\def\Ker{\mathop{\mathrm{Ker}}}
\def\Im{\mathop{\mathrm{Im}}}
\def\Coker{\mathop{\mathrm{Coker}}}
\def\Coim{\mathop{\mathrm{Coim}}}

% Boxtimes
%
\DeclareMathSymbol{\boxtimes}{\mathbin}{AMSa}{"02}

%
% Macros for moduli stacks/spaces
%
\def\QCohstack{\mathcal{QC}\!\mathit{oh}}
\def\Cohstack{\mathcal{C}\!\mathit{oh}}
\def\Spacesstack{\mathcal{S}\!\mathit{paces}}
\def\Quotfunctor{\mathrm{Quot}}
\def\Hilbfunctor{\mathrm{Hilb}}
\def\Curvesstack{\mathcal{C}\!\mathit{urves}}
\def\Polarizedstack{\mathcal{P}\!\mathit{olarized}}
\def\Complexesstack{\mathcal{C}\!\mathit{omplexes}}
% \Pic is the operator that assigns to X its picard group, usage \Pic(X)
% \Picardstack_{X/B} denotes the Picard stack of X over B
% \Picardfunctor_{X/B} denotes the Picard functor of X over B
\def\Pic{\mathop{\mathrm{Pic}}\nolimits}
\def\Picardstack{\mathcal{P}\!\mathit{ic}}
\def\Picardfunctor{\mathrm{Pic}}
\def\Deformationcategory{\mathcal{D}\!\mathit{ef}}
\usepackage{cleveref}
\crefname{originallemma}{original lemma}{original lemmas}
\Crefname{originallemma}{Original lemma}{Original lemmas}
\makeatletter
\newlabel{morphisms-lemma-Jacobson-universally-Jacobson}{{02J5}{source}{}{sourcepointer-1}{}}
\newlabel{morphisms-lemma-Jacobson-universally-Jacobson@cref}{{[originallemma][0][]02J5}{[1][0][]source}{}{}{}}
\makeatother
\crefname{originallemma}{original lemma}{original lemmas}
\Crefname{originallemma}{Original lemma}{Original lemmas}
\makeatletter
\newlabel{morphisms-lemma-jacobson-finite-type-points}{{01TB}{source}{}{sourcepointer-2}{}}
\newlabel{morphisms-lemma-jacobson-finite-type-points@cref}{{[originallemma][0][]01TB}{[1][0][]source}{}{}{}}
\makeatother
\crefname{originallemma}{original lemma}{original lemmas}
\Crefname{originallemma}{Original lemma}{Original lemmas}
\makeatletter
\newlabel{properties-lemma-codimension-local-ring}{{02IZ}{source}{}{sourcepointer-3}{}}
\newlabel{properties-lemma-codimension-local-ring@cref}{{[originallemma][0][]02IZ}{[1][0][]source}{}{}{}}
\makeatother
\crefname{originalreference}{original reference}{original references}
\Crefname{originalreference}{Original reference}{Original references}
\makeatletter
\newlabel{situation-setup}{{02QL}{source}{}{sourcepointer-4}{}}
\newlabel{situation-setup@cref}{{[originalreference][0][]02QL}{[1][0][]source}{}{}{}}
\makeatother
\title{Over a Jacobson base the relative dimension function agrees with the dimension of an integral subscheme}
\author{The Stacks Project Authors}
\date{stacks-project:chow.tex:local-snapshot-20261012; GNU FDL 1.2 or later}
\begin{document}
\maketitle
\noindent\textit{Editorial note.} All delivered mathematical text is the author's own, copied from the pinned Stacks Project chapter file chow.tex; no statement or proof step was rewritten, reordered or removed.\par\medskip
\noindent\textit{Editorial note.} Preamble provenance: the pinned chapter file opens with \textbackslash{}input\{preamble\} and carries no preamble text of its own, so the project's single Stacks-wide preamble staged at sources/stacks-chow/preamble.tex is excerpted verbatim as three spans, lines 8-12 (the verbatim package together with the project's own reference, slogan and history display-environment definitions, including the author's commented-out older form of the reference environment), lines 14-18 (Xy-pic setup) and lines 154-217 (the native theorem declarations and mathematical macros). The lines 8-12 span is required because a counted statement or its surrounding text may use those display environments, which the project typesets as comments. Nothing was added, removed or redefined; the wrapper class and the amsmath/amssymb/amsthm/graphicx base packages are supplied by the assembler.\par\medskip
\noindent\textit{Editorial note.} Same-chapter pointers are located at their \textbackslash{}label lines inside the pinned chapter file; cross-chapter pointers are located in the project's own tag list staged in the snapshot (sources/stacks-chow/tags.txt), at the verbatim lines carrying their canonical tags.\par\medskip
\noindent\textit{Allowed standard prerequisites (editorial).}\begin{itemize}
\item A scheme locally of finite type over a Jacobson scheme is Jacobson. Morphisms, label morphisms-lemma-Jacobson-universally-Jacobson, invoked at proof line 1475. Gives that closed points are dense in every closed subset and map to closed points of the base, so that a chain of irreducible closed subsets may be ended with a closed point.
\item The dimension of a scheme through the dimensions of its local rings, and finiteness of the residue field extension at a closed point. Properties, label properties-lemma-codimension-local-ring (proof line 1480); Morphisms, label morphisms-lemma-jacobson-finite-type-points (proof line 1486). The first expresses the dimension as the supremum of the dimensions of the local rings at the closed points, and the second makes the residue field extension at a closed point finite.
\end{itemize}\medskip
\ifdefined\hypertarget\hypertarget{sourcepointer-1}{}\fi
\noindent\textit{Source reference.} Original lemma 02J5: sources/stacks-chow/tags.txt, source lines 3324--3324. Navigation pointer for \textbackslash{}ref\{morphisms-lemma-Jacobson-universally-Jacobson\} used inside the counted span (line 1475); the argument that consumes it is retained in the span.\par\medskip
\ifdefined\hypertarget\hypertarget{sourcepointer-2}{}\fi
\noindent\textit{Source reference.} Original lemma 01TB: sources/stacks-chow/tags.txt, source lines 2368--2368. Navigation pointer for \textbackslash{}ref\{morphisms-lemma-jacobson-finite-type-points\} used inside the counted span (line 1486); the argument that consumes it is retained in the span.\par\medskip
\ifdefined\hypertarget\hypertarget{sourcepointer-3}{}\fi
\noindent\textit{Source reference.} Original lemma 02IZ: sources/stacks-chow/tags.txt, source lines 3318--3318. Navigation pointer for \textbackslash{}ref\{properties-lemma-codimension-local-ring\} used inside the counted span (line 1480); the argument that consumes it is retained in the span.\par\medskip
\ifdefined\hypertarget\hypertarget{sourcepointer-4}{}\fi
\noindent\textit{Source reference.} Original reference 02QL: sources/stacks-chow/pinned-original-source.tex, source lines 1372--1372. Navigation pointer for \textbackslash{}ref\{situation-setup\} used inside the counted span (line 1459); the argument that consumes it is retained in the span.\par\medskip
\begin{lemma}
\label{lemma-delta-is-dimension}
Let $(S, \delta)$ be as in Situation \ref{situation-setup}.
Assume in addition $S$ is a Jacobson scheme, and $\delta(s) = 0$ for every
closed point $s$ of $S$. Let $X$ be locally of finite type over $S$.
Let $Z \subset X$ be an integral closed subscheme and let
$\xi \in Z$ be its generic point. The following integers are the same:
\begin{enumerate}
\item $\delta_{X/S}(\xi)$,
\item $\dim(Z)$, and
\item $\dim(\mathcal{O}_{Z, z})$ where $z$ is a closed point of $Z$.
\end{enumerate}
\end{lemma}

\begin{proof}
Let $X \to S$, $\xi \in Z \subset X$ be as in the lemma.
Since $X$ is locally of finite type over $S$ we see that
$X$ is Jacobson, see
Morphisms, Lemma \ref{morphisms-lemma-Jacobson-universally-Jacobson}.
Hence closed points of $X$ are dense in every closed subset of $Z$
and map to closed points of $S$. Hence given any chain
of irreducible closed subsets of $Z$ we can end it with a closed point of $Z$.
It follows that $\dim(Z) = \sup_z(\dim(\mathcal{O}_{Z, z})$
(see Properties, Lemma \ref{properties-lemma-codimension-local-ring})
where $z \in Z$ runs over the closed points of $Z$.
Note that $\dim(\mathcal{O}_{Z, z}) = \delta(\xi) - \delta(z)$
by the properties of a dimension function.
For each closed $z \in Z$ the field extension
$\kappa(z)/\kappa(f(z))$ is finite, see Morphisms,
Lemma \ref{morphisms-lemma-jacobson-finite-type-points}.
Hence $\delta_{X/S}(z) = \delta(f(z)) = 0$ for $z \in Z$ closed.
It follows that all three integers are equal.
\end{proof}

\end{document}
